% ============================================================
% MÓDULO 4: PROGRAMACIÓN DINÁMICA Y GPI
% ============================================================
\section{M4 · Prog. Dinámica}
\module{Módulo 4: Programación Dinámica y GPI}{Cálculo de la política óptima cuando conocemos las leyes completas del mundo}

\begin{frame}{Ecuaciones de Bellman de Optimalidad}
  \begin{columns}[c,onlytextwidth]
    \column{0.50\textwidth}
      \begin{block}{Definición de Política Óptima $\pi^*$}
        Una política $\pi^*$ es óptima si y sólo si $V^{\pi^*}(s) \ge V^\pi(s)$ para todo estado $s$ y toda política $\pi$.
      \end{block}
      \vspace{0.2cm}
      Las funciones de valor óptimas satisfacen el principio de optimalidad:
      \[
        V^*(s) = \max_{a \in \mathcal{A}} Q^*(s,a)
      \]
    \column{0.48\textwidth}
      \begin{beamercolorbox}[rounded=true,sep=0.8em]{block body alerted}
        \centering
        {\large $V^*(s) = \displaystyle\max_{a} \left[ \mathcal{R}(s,a) + \gamma \sum_{s'} \mathcal{P}_{ss'}^a V^*(s') \right]$}
      \end{beamercolorbox}
      \vspace{0.25cm}
      \small \textbf{¡Ojo con la no linealidad!}\\
      Debido al operador $\max$, este ya \textbf{no} es un sistema de ecuaciones lineales. No se resuelve con inversión de matrices directa. Requiere métodos iterativos.
  \end{columns}
\end{frame}

\begin{frame}{Policy Iteration vs. Value Iteration}
  \begin{columns}[T,onlytextwidth]
    \column{0.48\textwidth}
      \begin{block}{Policy Iteration (Howard, 1960)}
        Alterna dos fases explícitas hasta converger:
        \begin{enumerate}
          \item \textbf{Evaluación de Política:} Resolver $V^{\pi_k}$ hasta convergencia.
          \item \textbf{Mejora de Política:}
          \[
            \pi_{k+1}(s) \leftarrow \arg\max_a Q^{\pi_k}(s,a)
          \]
        \end{enumerate}
        \textbf{Ventaja:} Pocas iteraciones de mejora.\\
        \textbf{Desventaja:} Evaluar $V$ en cada paso es costoso.
      \end{block}
    \column{0.48\textwidth}
      \begin{alertblock}{Value Iteration (Bellman, 1957)}
        Combina evaluación y mejora en un solo paso:
        \[
          V_{k+1}(s) \leftarrow \max_a \left[ \mathcal{R}(s,a) + \gamma \sum_{s'} \mathcal{P}_{ss'}^a V_k(s') \right]
        \]
        \textbf{Garantía matemática:} El operador de Bellman es una \textbf{contracción} en norma infinito ($L_\infty$). Converge geométricamente a la única solución $V^*$.
      \end{alertblock}
  \end{columns}
\end{frame}

\begin{frame}{Interacción Política Generalizada (GPI)}
  \centering
  \begin{tikzpicture}[scale=0.95,transform shape]
    \draw[Teal,line width=2.5pt] (0,0) -- (7,3.5) node[right,font=\bfseries,text=Teal] {$v \approx v_\pi$};
    \draw[Coral,line width=2.5pt] (0,3.5) -- (7,0) node[right,font=\bfseries,text=Coral] {$\pi \approx \text{greedy}(v)$};

    % Trayectoria en zigzag de GPI
    \draw[-{Latex},line width=1.5pt,Ink] (0.8,0.4) -- (1.5,2.75) node[midway,left,font=\scriptsize] {evaluación};
    \draw[-{Latex},line width=1.5pt,Ink] (1.5,2.75) -- (2.6,1.3) node[midway,above right,font=\scriptsize] {mejora};
    \draw[-{Latex},line width=1.5pt,Ink] (2.6,1.3) -- (3.5,1.75);
    \draw[-{Latex},line width=1.5pt,Ink] (3.5,1.75) -- (4.8,1.1);

    \fill[Gold] (3.5,1.75) circle (0.15cm);
    \node[above=0.2cm of {(3.5,1.75)},font=\bfseries\small,text=Gold!80!Ink] {Punto óptimo $(\pi^*, v^*)$};
  \end{tikzpicture}
  \vspace{0.4cm}
  \keyidea{Casi todo el RL consiste en GPI: dos procesos compitiendo y colaborando que convergen al punto fijo óptimo.}
\end{frame}

\begin{frame}[fragile]{Implementación en Python: Value Iteration para OptiBot}
  \begin{columns}[T,onlytextwidth]
    \column{0.62\textwidth}
\begin{lstlisting}[style=rlcode]
import numpy as np

# Cuadricula simple 3x3 del almacen
n_states = 9
gamma = 0.95
V = np.zeros(n_states)
theta = 1e-4  # Umbral de convergencia

# Bucle de Value Iteration
while True:
    delta = 0
    for s in range(n_states):
        v = V[s]
        # Calcular valor para cada accion posible
        # Q(s, a) = R(s,a) + gamma * sum(P * V)
        q_values = []
        for a in range(4): # 0:Arriba, 1:Abajo, 2:Izq, 3:Der
            r = -1.0 if s != 8 else 0.0 # Recompensa
            s_next = s # Logica de transicion simple
            q_values.append(r + gamma * V[s_next])
        V[s] = max(q_values)
        delta = max(delta, abs(v - V[s]))
    if delta < theta:
        break
\end{lstlisting}
    \column{0.35\textwidth}
      \begin{alertblock}{La limitación de M4}
        DP es elegante y exacto.\\
        \vspace{0.2cm}
        \textbf{Pero exige conocer $\mathcal{P}$ y $\mathcal{R}$ de antemano.}\\
        \vspace{0.2cm}
        En un almacén real con obstáculos móviles y suelo irregular, \textbf{nadie} nos entrega la matriz de transición.\\
        \vspace{0.2cm}
        ¿Cómo aprende OptiBot sólo interactuando por ensayo y error? $\to$ \textbf{Monte Carlo (M5)}.
      \end{alertblock}
  \end{columns}
\end{frame}
